\section{Measurements, validation, and repository integration}
\label{sec:experiments}

The experiments separate three computational questions.  The window
benchmark measures the cost of obtaining a selected obstruction instead
of all homology.  The perturbation benchmark computes the same complete
homology by two reduction engines.  The pattern benchmark compares two
classifiers for an already known ball.  A speed ratio from one question
is not an end-to-end speed ratio for the other two.

All stored measurements use Python 3.12.14 on an x86-64 Linux environment
with glibc 2.39.  Raw repetitions, input diagrams or deterministic fixture
generators, crossing orders where applicable, and operation counts are
included.  Wall times are observations on this environment.  The paired
protocol reduces some variability but does not remove shared-machine
scheduling noise, especially for millisecond inputs.  Asymptotic claims
in this paper follow from the preceding proofs, rather than fits to these
small timing samples.

\subsection{Exact windows: less construction, sometimes an earlier verdict}

The window experiment uses one warmup per variant and seven timed pairs,
alternating which variant runs first.  Both variants receive identical
PD data and the same preselected crossing order.  Both use the baseline's
automatic shape-cache policy.  Diagram loading, order selection, and
cheap invariant filters are excluded.  The window orientation is fixed
for each row: the bound-aware automatic mirror selector is available to
callers but is not being credited for these measurements.

\begin{table}[htbp]
\centering\small
\setlength{\tabcolsep}{4pt}
\begin{tabular}{lrrrrr}
\toprule
Input & $n$ & Full (ms) & Window (ms) & Full/window & $\dim H^0$ \\
\midrule
Conway & 11 & 7.202 & 7.089 & 1.016 & 10 \\
Kinoshita--Terasaka & 11 & 8.954 & 7.747 & 1.156 & 10 \\
Hard unknot & 8 & 1.626 & 1.364 & 1.193 & 2 \\
Determinant-one grid & 10 & 4.826 & 1.455 & 3.316 & 2 \\
Mixed 5-braid & 36 & 904.254 & 581.562 & 1.555 & 698 \\
Conway, mirrored & 11 & 13.823 & 13.456 & 1.027 & 10 \\
K--T, mirrored & 11 & 13.184 & 14.254 & 0.925 & 10 \\
$T(4,7)$ control & 21 & 77.047 & 9.280 & 8.303 & 2 \\
$T(2,31)$ control & 31 & 11.580 & 6.467 & 1.791 & 2 \\
\bottomrule
\end{tabular}
\caption{Scanner-kernel medians from seven alternating pairs after a warmup of each variant. The window computes only normalized degree zero, while the full scan computes every degree. Input preparation and order selection are excluded.}
\label{tab:window-times}
\end{table}

\begin{table}[htbp]
\centering\small
\setlength{\tabcolsep}{4pt}
\begin{tabular}{lrrc}
\toprule
Input & Full peak & Window peak & Window rejects unknot? \\
\midrule
Conway & 246 & 124 & Yes \\
Kinoshita--Terasaka & 246 & 118 & Yes \\
Hard unknot & 11 & 11 & No \\
Determinant-one grid & 54 & 24 & No \\
Mixed 5-braid & 17,694 & 7,692 & Yes \\
Conway, mirrored & 246 & 128 & Yes \\
K--T, mirrored & 246 & 122 & Yes \\
$T(4,7)$ control & 210 & 40 & No \\
$T(2,31)$ control & 186 & 66 & No \\
\bottomrule
\end{tabular}
\caption{Peak physical objects allocated before elimination. These are object counts, not peak resident-memory measurements. A degree-zero dimension of two is inconclusive.}
\label{tab:window-objects}
\end{table}


The most useful case is the repository's 36-crossing mixed five-braid
stress diagram.  Its complete unreduced rank is 5898, while its normalized
degree-zero rank is 698.  Thus this one degree already proves
nontriviality.  The median full scan took $0.904254$ seconds and the
window took $0.581562$ seconds, a ratio of approximately $1.55$.
The peak allocated objects decreased from 17694 to 7692, a reduction of
about $56.5\%$.  Figure~\ref{fig:window-profile} plots every allocation
stage, including the large final expansion.

The identical-cap experiment is more informative than a small timing
difference.  With a 10000-object ceiling the full scan raises
\texttt{ScanLimit} before allocating 17694 objects, while the window
completes with peak 7692 and returns the 698-dimensional obstruction.
This is a reproducible change in which exact question can be answered
within a stated resource limit.  The cap is an object cap, not a measured
resident-memory limit; morphisms, dictionaries and caches also use space.

Conway and Kinoshita--Terasaka both have degree-zero dimension 10 in these
conventions, so their window results are also decisive.  Their short wall
times should not be overinterpreted: the mirrored Kinoshita--Terasaka row
is slower with the window.  Truncation incurs bookkeeping, and different
remaining objects can change the pivot order.

The hard unknot, determinant-one grid knot, and torus controls all have
degree-zero dimension 2.  Those window results are inconclusive.  In
particular, the large ratio on the $T(4,7)$ control does not constitute
a faster complete recognition result: additional homology or a cheap
invariant still supplies nontriviality.  The standard recognition pipeline
already dispatches many of these inputs before Khovanov computation.
These rows deliberately isolate the scanner operation.

The saved JSON also contains the isolated-guard pruning ablation.  Every
ablation agrees on the requested ranks.  Its allocation profile explains
what this additional operation deletes, but the ablation was not timed
and no speed ratio is assigned to it.  For the final stress stage the
retained allocation is 7692; unit elimination and the valid deletion of
isolated guards leave the requested 698 objects.

\subsection{Perturbation: an exact prototype with a measured tradeoff}

The second experiment uses five warmed alternating pairs.  The complete
scanner and perturbation backend receive the same fixed crossing order,
with shape caching disabled in both.  They agree on the full homological
rank distribution for every row.  Optional construction and verification
of all contraction maps is excluded from the timing and tested separately.

\begin{table}[htbp]
\centering\small
\setlength{\tabcolsep}{4pt}
\begin{tabular}{lrrrr}
\toprule
Input & $n$ & Scalar (ms) & Perturbation (ms) & Scalar/pert. \\
\midrule
Hard unknot & 8 & 2.299 & 3.479 & 0.661 \\
Conway & 11 & 7.858 & 15.080 & 0.521 \\
Kinoshita--Terasaka & 11 & 8.572 & 19.468 & 0.440 \\
$T(3,7)$ & 14 & 13.637 & 15.951 & 0.855 \\
$T(4,7)$ & 21 & 45.236 & 60.884 & 0.743 \\
\bottomrule
\end{tabular}
\caption{Full Khovanov scans: five warmed alternating pairs, fixed order, shape cache disabled in both backends. The perturbation backend is slower on every natural input tested.}
\label{tab:hpl-natural}
\end{table}


The current perturbation implementation loses on every natural diagram
in Table~\ref{tab:hpl-natural}.  The extra binary basis construction and
map assembly outweigh the avoided individual cobordism operations in
these cases.  This is a reason to retain it as an explicit research
backend, and to use its rank profiles and certificates to design a
selective hybrid, rather than replace the default elimination routine.

There is nevertheless a regime where the construction has a concrete
advantage.  Consider two-term algebraic complexes with differential
\[
 d=I+xR\quad\text{over }\mathbb F_2[x]/(x^2),
\]
where $R$ is a deterministic pseudorandom binary matrix.  The scalar
differential is the identity, so its homology vanishes.  Moreover
$(I+xR)^2=I$, making full contractibility immediate.  These are valid
complexes over the one-arc endomorphism algebra, although no claim that
they arise from knot diagrams is made.

\begin{table}[htbp]
\centering\small
\setlength{\tabcolsep}{4pt}
\begin{tabular}{rrrrr}
\toprule
$R$ order & Objects & Scalar (ms) & Perturbation (ms) & Scalar/pert. \\
\midrule
16 & 32 & 0.162 & 0.192 & 0.842 \\
32 & 64 & 0.744 & 0.444 & 1.676 \\
64 & 128 & 4.022 & 1.449 & 2.776 \\
128 & 256 & 29.020 & 5.822 & 4.984 \\
256 & 512 & 204.609 & 31.587 & 6.478 \\
\bottomrule
\end{tabular}
\caption{Synthetic two-term complexes with differential $I+xR$ over $\mathbb F_2[x]/(x^2)$. These are algebraic stress tests, not knot diagrams. Both timings include construction of the same input complex.}
\label{tab:hpl-synthetic}
\end{table}


At matrix order 256, the median decreases from $0.204609$ to $0.031587$
seconds, about $6.48$ times.  Both timings construct the same input.
The scalar path performs many nilpotent updates; the perturbation path
detects scalar acyclicity and returns an empty minimal object.  Requiring
an explicit contracting homotopy would add work.  The benefit therefore
supports a specific implementation hypothesis---large scalar rank can
make block reduction attractive---without supplying an end-to-end
unknot speedup.  The stored small stress cases do check the complete
contracting identities.

\subsection{Terminal patterns: a large improvement in a narrow subroutine}

\begin{table}[htbp]
\centering\small
\setlength{\tabcolsep}{4pt}
\begin{tabular}{rrrr}
\toprule
Darts & New (ms) & Report 06 (ms) & Old/new \\
\midrule
24 & 0.062 & 2.526 & 40.6 \\
48 & 0.081 & 25.605 & 315.9 \\
72 & 0.108 & 152.797 & 1416.0 \\
96 & 0.149 & 487.551 & 3265.9 \\
144 & 0.166 & 2745.396 & 16517.9 \\
192 & 0.440 & 9530.512 & 21650.6 \\
\bottomrule
\end{tabular}
\caption{Known-ball terminal classifier: medians of three complete calls including validation. The fixture family is dual to bipyramid triangulations. No complete hierarchy is timed.}
\label{tab:terminal-times}
\end{table}

\begin{table}[htbp]
\centering\small
\setlength{\tabcolsep}{4pt}
\begin{tabular}{rrr}
\toprule
Darts & New (ms) & Candidate pairs \\
\midrule
384 & 0.526 & 193 \\
1536 & 5.704 & 769 \\
6144 & 23.201 & 3073 \\
24576 & 133.726 & 12289 \\
98304 & 496.136 & 49153 \\
\bottomrule
\end{tabular}
\caption{Larger essential terminal fixtures, without running the quartic baseline. Counts and individual timings are in the raw JSON.}
\label{tab:terminal-scaling}
\end{table}


The old and new pattern timings both begin with the same prepared
JSON-like input and include validation, spherical-embedding checks,
classification, and result construction.  Fixture generation and
external JSON reading and writing are excluded.  The baseline source is
the unmodified report-06 implementation, preserved with Git blob hashes.
Each median uses three calls.

At 192 darts, the new classifier takes about $0.440$ milliseconds and the
old bond enumerator about $9.53$ seconds.  The very large ratio is
consistent with eliminating an enumeration of edge subsets followed by
repeated graph traversal.  On the larger fixture with 98304 darts the new
classifier takes approximately $0.50$ seconds.  These fixtures provide
useful checks of the linear-space and near-linear-time construction.
They give no estimate for normal-surface selection, hierarchy length, or
the number of terminal components in a complete recognition algorithm.

\subsection{Validation and the remaining trust boundary}

The delivered fast-scanner suite passes 64 test methods, comprising the
41 baseline methods and 23 additions.  The additions cover exact windows,
mirror selection and normalization, guarded boundaries, truncation before
resource accounting, randomized scan orders, cache choices, full-fallback
semantics, the CLI, and finite perturbation identities.  The hierarchy
suite adds 10 passing methods.  The full logs are included; individual
methods may contain many exhaustive or randomized cases.

Two especially relevant checks are more discriminating than agreement on
one final integer.  First, intermediate perturbation models have the same
matching-by-degree multiplicities as exhaustive unit cancellation, and
their explicit maps satisfy all deformation-retraction identities.
Second, all 10410 small rotation maps are censused by the terminal tests,
and connected negative witnesses are verified by the old primal-cut
checker as well as the new verifier.  The tests also reject deliberately
corrupted certificates and nonspherical inputs.

The saved HPL example is the unreduced four-crossing prefix of the
eight-crossing $(3,[1,2]^4)$ braid closure.  It reduces 20 source objects
to 6 minimal objects at a six-endpoint frontier, and its constructor
uses perturbation depth two.  Its JSON records the source diagram,
actual endpoint pairings, and all five differential and contraction
maps.  The verifier reconstructs the source prefix and checks the saved
maps without recomputing the perturbation construction.  A second JSON
records a terminal triangle with primal sides of sizes three and three.
Both artifacts can be reloaded and checked with
\path{experiments/export_certificates.py --verify-existing}.

The proofs in this article are mathematical arguments accompanied by
executable checks.  They are not Lean formalizations.  The window API
returns ranks, a degree convention, and allocation evidence; it does not
yet export a short, standalone proof of those ranks checkable without the
scanner.  The perturbation certificate checks trust the exact implemented
cobordism composition.  The terminal negative verifier is independent of
the triangle-enumeration algorithm but still trusts its input's known-ball
precondition.  These are specific interfaces for further formalization,
not substitutes for the hypotheses in the correctness theorems.

\subsection{Integration contract}

The complete modified \texttt{fast/} directory is included, together with
\path{integration/fast-window.patch}.  The patch is based on the active
tree at the pinned revision, and its applicability is checked against that
exact snapshot.  The report-07 structural patch is a separate lineage;
combining both continuations requires an ordinary code integration and
the corresponding combined regression run.  No remote repository is
changed by preparing this package.

The new recognition stage is opt-in.  Existing calls preserve their
default behavior.  The main API additions are:
\begin{lstlisting}[language=Python]
recognize(diagram, use_window=True, window_radius=0,
          window_mirror=True, window_max_objects=10000,
          window_seconds=None)
\end{lstlisting}
The target is the normalized interval
$[-\texttt{window\_radius},\texttt{window\_radius}]$.
Automatic mirroring compares the exact proved bounds $B_0,B_1$ and
restores the original raw indices before normalization.  A result
different from $\{0:2\}$ is a nontriviality obstruction.  Agreement invokes
the full fallback.  Exhausting a window-specific budget records a skipped
window and leaves the remaining global budget to that fallback.  Exhausting
the complete budget yields \texttt{UNKNOWN}.

The lower-level routines distinguish raw-degree queries from this
normalized decision interface:
\begin{lstlisting}[language=Python]
from fastunknot.window_scan import khovanov_window
from fastunknot.window_bounds import khovanov_window_auto

# Raw degrees, fixed orientation:
part = khovanov_window(diagram.pd, lower=a, upper=b,
                       order=order, max_objects=10000)

# The same original raw degrees, selecting a mirror internally:
part = khovanov_window_auto(diagram, a, b, order=order)
\end{lstlisting}
The result field is \texttt{window\_rank}.  A separate
\texttt{complete\_rank} flag is true only when the requested interval
covers the entire possible raw range $[0,n]$.  The \texttt{window} CLI
subcommand also accepts raw bounds, recorded as \texttt{raw\_lower} and
\texttt{raw\_upper}; it does
not label a partial answer as a total rank.  The complete recognizer's
\texttt{--window} flag enables the optional obstruction stage.

The module \texttt{perturbation.py} is available for direct research use
and certified intermediate reductions; it is not selected by the default
recognizer.  The \texttt{hierarchy/} directory is a self-contained terminal
classifier and experiment bundle.  It is intended for a future hierarchy
implementation or a report appendix with its ambient-ball contract intact.

\paragraph{Reproducibility artifacts.}
The package includes the TeX source and figure PDFs, raw timing JSON,
deterministic fixture generators, source provenance, the integration patch,
and complete tests.  The \texttt{experiments/} scripts regenerate benchmark
data, figures, tables and example certificates.  Figure generation checks
the full allocation trace against the saved stress rank and peak before
plotting it.  The tables are generated directly from the raw JSON so that
the numerical claims can be updated without manual transcription.

These deliverables support two immediate engineering choices: enable the
window obstruction for workloads where it can avoid an expensive full
scan, and replace repeated terminal bond enumeration where the known-ball
precondition is already certified.  The nilpotent backend supplies a
separate algebraic foundation for the next optimization cycle.
